%% ──────────────────────────────────────────────────────────────────
%%  Diabetic Retinopathy — Lab 5 Report (v3)
%%  Computer Vision, April 2026
%% ──────────────────────────────────────────────────────────────────
\documentclass[a4paper, 10pt, twocolumn]{article}

\usepackage[top=1.6cm, bottom=1.6cm, left=1.4cm, right=1.4cm, columnsep=0.65cm]{geometry}
\usepackage[utf8]{inputenc}
\usepackage[T1]{fontenc}
\usepackage{lmodern}
\usepackage{microtype}
\usepackage{booktabs}
\usepackage{multirow}
\usepackage[table]{xcolor}
\usepackage{titlesec}
\usepackage[hidelinks]{hyperref}
\usepackage{enumitem}
\usepackage{amsmath}
\usepackage[labelfont=bf, font=small, skip=3pt]{caption}
\usepackage{tabularx}
\usepackage{array}
\usepackage{graphicx}

%% ── Colour palette ────────────────────────────────────────────────
\definecolor{accent}{RGB}{0, 64, 128}
\definecolor{rowbest}{RGB}{220, 240, 220}
\definecolor{rowbad} {RGB}{255, 230, 230}
\definecolor{rownorm}{RGB}{250, 250, 252}
\definecolor{rulecolor}{RGB}{0, 64, 128}

%% ── Section headings ─────────────────────────────────────────────
\titleformat{\section}
  {\normalfont\bfseries\small\color{accent}}
  {}{0em}{}[{\color{rulecolor}\vspace{-1pt}\hrule height 0.6pt\vspace{1pt}}]
\titlespacing{\section}{0pt}{4pt}{1pt}

\titleformat{\subsection}
  {\normalfont\small\bfseries}
  {}{0em}{}
\titlespacing{\subsection}{0pt}{2pt}{0pt}

%% ── Paragraph / list style ───────────────────────────────────────
\setlength{\parskip}{1pt}
\setlength{\parindent}{0pt}
\setlist[itemize]{topsep=0pt, itemsep=0pt, leftmargin=10pt,
                  label=\raisebox{0.3ex}{\tiny$\bullet$}, parsep=0pt}

%% ── Captions ─────────────────────────────────────────────────────
\setlength{\abovecaptionskip}{3pt}
\setlength{\belowcaptionskip}{0pt}

%% ══════════════════════════════════════════════════════════════════
\begin{document}

%% ── Title spanning both columns 
\twocolumn[{%
\begin{center}
  {\large\bfseries\color{accent}
    LS5: Image Classification with CNNs \\[2pt] Automatic Diagnosis System for Diabetic Retinopathy}\\[4pt]
  {\small
    Maria Muñoz Pérez (100471836) -
    Cristina Sobrino Fernández-Miranda (100471865)}\\[3pt]
  {\color{rulecolor}\hrule height 0.8pt}
  \vspace{5pt}
\end{center}
}]

%% ── 1. INTRODUCTION & DATA ──────────────────────────────────────
\section{Introduction}

We address binary classification of Diabetic Retinopathy from
fundus photographs (2\,000 train / 500 val / 1\,000 test, 73\%/27\%
class imbalance) under two competition tracks: a CNN trained from scratch
(\textbf{CUSTOM}) and a fine-tuned ImageNet backbone (\textbf{FINE-TUNING}).
Below we describe our final solution for each track and a summary of the failed approaches.

%% ── 2. CUSTOM MODEL ────────────────────────────────────────────
\section{CUSTOM: 3-Seed CustomNetV2 Ensemble}


\textbf{Preprocessing.}
Each image is cropped to the retinal disc via intensity thresholding
(\texttt{CropByEye}), then contrast-enhanced with the Ben Graham
formula~\cite{graham2015kaggle}
($4{\times}\text{img}-4{\times}\text{GaussianBlur}+128$, with circular
mask and a \texttt{uint8} dtype guard). The result is rescaled to 256\,px, followed by random
horizontal and vertical flips, $\pm$15° rotation, colour jitter, and a
random 224\,px crop at training (centre crop at inference).
\texttt{WeightedRandomSampler} and
\texttt{BCEWithLogitsLoss(pos\_weight$\approx$2.7)} together handle
the class imbalance.

\textbf{Architecture.}
Five sequencial blocks. Each block includes a $3\times3$ convolution (with same padding and no bias), Batch Normalization, ReLU activation, and $2\times2$ Max Pooling, with channel depth expanding from
$3{\to}32{\to}64{\to}128{\to}256{\to}256$. We concluded the network with a Global Average Pooling (GAP) layer, which flattens the spatial output, followed by a classification head consisting of two fully connected layers with dropout ($p=0.2$ and $p=0.4$) to try to mitigate overfitting (\textbf{Total: $1,012,257$ parameters}).
Training uses AdamW (lr=1e\textminus3, wd=5e\textminus3) with
CosineAnnealingLR($T_\text{max}$=50), 50 epochs and a patience of 10.

\textbf{Ensemble.}
Three independent runs (seeds 42, 123, 456) are averaged, each with
4-pass TTA (original, hflip, vflip, rot180). 
The diversity resulting from independent weight initializations effectively mitigates the risk of  overfitting. We rely solely on a fixed uniform average, avoiding any performance-based model selection.


%% ── 3. FINE-TUNING MODEL ───────────────────────────────────────
\section{FINE-TUNING: 3-Backbone Weighted Ensemble}

\textbf{Backbones:} Three ImageNet-pretrained backbones with complementary
inductive biases are combined: 
\begin{enumerate}[topsep=0pt, itemsep=1pt]
    \item \textbf{DenseNet121}~\cite{huang2017densely} (dense feature reuse),
    \item  \textbf{DenseNet169} (deeper dense reuse)
    \item \textbf{ResNeXt50-32x4d} (grouped convolutions). 
\end{enumerate}


Each has a custom head \texttt{Dropout(0.5) $\to$ Linear($d$,1)} where $d$ is the backbone's feature dimension (1024/1664/2048).
    
\textbf{Two-stage progressive unfreezing:} We train each model in two stages. Stage 1 freezes the backbone except for the classification head and the deepest block (AdamW, lr=3e-4, 30 epochs, label smoothing $\varepsilon=0.05$). Stage 2 unfreezes one additional block, reducing the backbone learning rate to 3e-5 and the head to 1e-4. We incorporate a 3-epoch linear warm-up to prevent gradient spikes upon unfreezing. To ensure stability, a validation gate skips Stage 2 if Stage 1 performance is insufficient (val AUC < 0.752), and we restore the Stage 1 checkpoint if Stage 2 regresses.

\textbf{High-resolution pipeline (380\,px).} The primary driver of sensitivity for early DR
lesions are microaneursysms~\cite{abramoff2018pivotal}, which are often lost at lower resolutions, so we scale inputs to 380px. Our pipeline uses: \texttt{CropByEye} $\to$
\texttt{BenGraham} $\to$ \texttt{Rescale(416)} $\to$ random hflip $\to$
\textbf{rotation\,$\pm$180$^\circ$} $\to$ ColorJitter $\to$
\texttt{RandomCrop(380)} $\to$ \texttt{Normalize}. The full $\pm$180$^\circ$
range is used because fundus cameras have no canonical orientation, the
standard $\pm$15$^\circ$ window discarded valid rotations,
and full rotation is reported as the an impactful fundus
augmentation~\cite{graham2015kaggle}. Backbones are fully convolutional
(GAP head), so no architectural change was needed for the higher resolution.\\
\textbf{Inference \& ensemble:} Predictions are processed with 4-pass TTA. Using a grid search on the 500-image validation set, the optimal weighted ensemble is:$$0.35 \times \text{DenseNet121} + 0.15 \times \text{DenseNet169} + 0.50 \times \text{ResNeXt50}.$$
We explicitly excluded another trained backbone (SE-ResNeXt50), as its high correlation with ResNeXt50 reduced the ensemble's diversity. Val AUC 0.867. 

\begin{table}[htbp]
    \centering
    \label{tab:final_results}
    \begin{tabular}{lrr}
        \toprule
        \textbf{Track} & \textbf{Public AUC} & \textbf{Private AUC} \\
        \midrule
        CUSTOM & 0.757 &\textbf{ 0.780} \\
        FINE-TUNING    & 0.835 & \textbf{0.826} \\
        \bottomrule
    \end{tabular}
\end{table}

\section{What Did Not Work}

The following approaches were tested and discarded (see Table\,\ref{tab:abl}):
\begin{itemize}
  \item \textbf{CLAHE}: altered intensity distribution, broke \texttt{CropByEye} crops.
  \item \textbf{BenGraham uint8 bug}: pixel overflow produced uniform grey images (custom 0.541, FT 0.558). Fixed with a dtype guard.
  \item \textbf{Pre-cached 224\,px images}: removed per-epoch \texttt{RandomCrop}, the main spatial regulariser ($-$0.028 AUC).
  \item \textbf{SE blocks / residual connections (custom)}: overfitted on 2\,000 images without ImageNet pretraining; SE-V2 reached only 0.52 val AUC.
  \item \textbf{380\,px + $\pm$180° on custom model}: failed to converge (3-seed val AUC 0.53); aggressive augmentation only helps pretrained backbones.
  \item \textbf{8-pass TTA (rot90/rot270)}: 90° rotations fall outside the training distribution due to BenGraham + rectangular crops; AUC dropped on both tracks.
  \item \textbf{DenseNet Stage\,3 unfreeze}: overfitted on 2\,000 images ($-$0.013 FT AUC); two stages are the right stopping point.
\end{itemize}

%% ══════════════════════════════════════════════════════════════════
%% PAGE 2 — TABLES AND REFERENCES
%% ══════════════════════════════════════════════════════════════════
\newpage
\onecolumn

\section*{\color{accent}Tables and References}
{\color{rulecolor}\hrule height 0.6pt}
\vspace{4pt}

%% ── TABLE 1: Results progression ──────────────────────────────────
\captionof{table}{Key Codabench submissions showing the AUC trajectory
  (\emph{pub} = public leaderboard, \emph{priv} = private).
  Final solutions highlighted in green.
  $^\dagger$\,Custom result invalid: CSV blended pretrained (EfficientNet-B0) scores.}
\label{tab:traj}
\vspace{3pt}
\footnotesize
\setlength{\tabcolsep}{5pt}
\renewcommand{\arraystretch}{1.12}
\begin{tabularx}{\textwidth}{@{}l X cc@{}}
\toprule
\textbf{Notebook} & \textbf{Key change} &
  \textbf{Custom AUC} & \textbf{FT AUC} \\
 & & \scriptsize pub\;/\;priv & \scriptsize pub\;/\;priv \\
\midrule
\texttt{first\_attempt}
  & CustomNet baseline + EfficientNet-B0 (SGD + StepLR)
  & 0.570\,/\,0.597 & 0.698\,/\,0.708 \\
\texttt{second\_attempt\_tuned}
  & CustomNetV2 + BenGraham + WeightedRandomSampler + 4-pass TTA
  & 0.766\,/\,0.743 & 0.770\,/\,0.744 \\
\texttt{fifth\_attempt}
  & BenGraham dtype guard fixed; label smoothing added (FT)
  & \textit{0.770\,/\,0.770}$^\dagger$ & 0.757\,/\,0.755 \\
\texttt{eighth\_attempt}
  & Custom: SE-V2+V3 ensemble;\enspace FT: ResNet50 two-stage unfreezing
  & 0.756\,/\,0.764 & 0.780\,/\,0.791 \\
\texttt{ninth\_attempt}
  & Custom: 3-seed V2;\enspace FT: 5-backbone 380\,px $\pm$180° (attempt\_23)
  & 0.757\,/\,0.780 & 0.806\,/\,0.807 \\
\rowcolor{rowbest}
\texttt{ninth\_attempt}
  & \textbf{CUSTOM final} --- 3-seed CustomNetV2, best private AUC
  & \textbf{0.757\,/\,0.780} & --- \\
\rowcolor{rowbest}
\texttt{attempt\_24}
  & \textbf{FT final} --- val-grid 3-backbone weighted ensemble (380\,px)
  & --- & \textbf{0.835\,/\,0.826} \\
\bottomrule
\end{tabularx}

\vspace{10pt}

%% ── TABLE 2: Ablations ─────────────────────────────────────────────
\captionof{table}{Failed approaches discussed in \S\,4.
  \textit{Cust} = CUSTOM, \textit{FT} = FINE-TUNING, \textit{Both} = both tracks.}
\label{tab:abl}
\vspace{3pt}
\footnotesize
\setlength{\tabcolsep}{4pt}
\renewcommand{\arraystretch}{1.18}
\begin{tabularx}{\textwidth}{@{}p{0.85cm} p{5.4cm} X@{}}
\toprule
\textbf{Track} & \textbf{Change} & \textbf{Why discarded} \\
\midrule

Both
  & CLAHE before \texttt{CropByEye}
  & Alters global intensity; \texttt{CropByEye}'s fixed threshold (0.10)
    mis-segments the retinal disc. \\

Both
  & BenGraham without \texttt{uint8} dtype guard
  & \texttt{uint8} overflow collapses all pixels to $\{0,\,128\}$ via silent
    numpy arithmetic, producing uniform grey images (custom 0.541, FT 0.558). \\

Both
  & Pre-cached 224\,px images on disk (5$\times$ speed-up)
  & \texttt{CenterCrop} at cache time removes per-epoch \texttt{RandomCrop}
    jitter, the primary spatial regulariser on a 2\,000-image dataset ($-$0.028 AUC). \\

Cust
  & SE blocks (r\,=\,8) on all CustomNetV2 blocks + residual connections (CustomNetV3)
  & SE-V2 alone reached 0.52 val AUC; residuals gave at best $+$0.002.
    Extra parameters overfit on 2\,000 images without ImageNet pretraining. \\

Cust
  & 380\,px input + $\pm$180° rotation on CustomNetV2
  & The 5$\times$MaxPool path is tuned for a 7$\times$7 feature map at 224\,px;
    at 380\,px training from scratch diverges (3-seed val AUC 0.53). \\

Both
  & 8-pass TTA (rot90 / rot270 added to 4-pass set)
  & BenGraham + rectangular crops are not invariant under 90° rotation;
    augmented images fall outside the training distribution. \\

FT
  & DenseNet Stage\,3 unfreeze (\texttt{denseblock2 + transition2})
  & Over-parameterises on 2\,000 images; $-$0.013 FT AUC vs.\ the two-stage strategy. \\

\bottomrule
\end{tabularx}

\vspace{10pt}

%% ── REFERENCES ───────────────────────────────────────────────────
\renewcommand{\refname}{\normalfont\bfseries\small\color{accent}References}
\begin{thebibliography}{9}
\small\setlength{\itemsep}{1pt}

\bibitem{graham2015kaggle}
Graham, B.\,(2015).
Kaggle diabetic retinopathy detection competition report.
\textit{University of Warwick}.

\bibitem{huang2017densely}
Huang, G., Liu, Z., van der Maaten, L., \& Weinberger, K.\,Q.\,(2017).
Densely connected convolutional networks.
\textit{Proc.\ CVPR 2017}, 4700--4708.

\bibitem{abramoff2018pivotal}
Abramoff, M.\,D., et al.\,(2018).
Pivotal trial of an autonomous AI-based diagnostic system for detection of
diabetic retinopathy in primary care offices.
\textit{npj Digital Medicine}, 1(1), 39.

\bibitem{loshchilov2019adamw}
Loshchilov, I., \& Hutter, F.\,(2019).
Decoupled weight decay regularization.
\textit{Proc.\ ICLR 2019}.

\end{thebibliography}

\end{document}
